\documentclass[11pt]{article}
\usepackage[margin=1in]{geometry}
\usepackage{amsmath,amssymb,amsthm,booktabs,array}
\usepackage[T1]{fontenc}
\usepackage{lmodern}
\usepackage{microtype}
\usepackage[hidelinks]{hyperref}
\hypersetup{
pdftitle={Hardin's trace-avoiding array recurrences and Barker's formulas: alternative proofs extending Riedel's and Howroyd's enumerations},
pdfauthor={Haobo Ma; Wenlin Zhang},
pdfsubject={Trace-avoiding binary arrays; CC BY 4.0}}
\usepackage{fancyvrb}
\newtheorem{theorem}{Theorem}
\newtheorem{lemma}[theorem]{Lemma}
\newtheorem{corollary}[theorem]{Corollary}
\newcommand{\one}{\mathbf{1}}
\newcommand{\leanbase}{https://github.com/the-omega-institute/trureturing/blob/1a461fb83236b7d5f88b0ffd1bde5f7d1e3ff8b4/D5/S3/Combinatorics/TraceAvoidingArrays/}
\title{Hardin's trace-avoiding array recurrences and Barker's formulas:\\
alternative proofs extending Riedel's and Howroyd's enumerations}
\author{Haobo Ma\\\small ChronoAI Pte Ltd, Singapore
\and Wenlin Zhang\\\small National University of Singapore\\
\small ChronoAI Pte Ltd, Singapore}
\date{14 November 2026}
\begin{document}
\maketitle
\begin{abstract}
We give alternative proofs and extensions for R.\ H.\ Hardin's empirical recurrences in the formula fields of
OEIS A185761 (revision 8), A185762--A185763 (revision 9), A185764
(revision 10), and A185765--A185768 (revision 9), together with Colin
Barker's formulas added on 22 February and 16 April 2018; the width-two
case follows from Marko Riedel's 2016 word enumeration and Andrew
Howroyd's 2025 array argument.
These sequences count labeled binary arrays whose neighboring contiguous
$2\times2$ blocks have different ordinary traces.
A row-pair transfer construction reduces the assertions at widths two
through nine to explicit finite integer identities; multiplying those
identities by further transfer powers proves the recurrences at every
length. It also gives rational generating functions at every fixed width
and explains the transposition symmetry of the parent table A185769.
\end{abstract}

\section{The counts and the formulas}
Hardin's entries~\cite{oeis2,oeis3,oeis4,oeis5,oeis6,oeis7,oeis8,oeis9} give empirical recurrences for
trace-avoiding binary arrays, and Barker gives the corresponding generating
functions. We establish these formulas by a transfer argument with exact
finite certificates. The recurrence and generating-function expressions
are those in the entries, rather than newly proposed formulas.
The method is the usual finite-digraph transfer method
\cite[Theorems 4.7.1--4.7.2]{stanley}; the specific continuation identities
below establish the empirical formulas with their exact starting indices.
Riedel's dated word-counting argument~\cite{riedel} already implies the
width-two recurrence and generating function, and hence its closed form,
through the explicit encoding below. Howroyd's later array-counting
argument~\cite{howroyd} gives another antecedent for that case.
The present uniform construction
extends that enumeration to widths three through nine; the corresponding
fixed-revision Lean statements for widths two through six are identified
in Section~4.

For integers $w\ge2$ and $n\ge1$, let $a_w(n)$ be the number of arrays
$A=(A_{ij})$ with $n+1$ rows, $w$ columns, and entries in $\{0,1\}$, for which
\[
 t_{ij}=A_{ij}+A_{i+1,j+1}\qquad(0\le i<n,\quad 0\le j<w-1)
\]
satisfies $t_{ij}\ne t_{i,j+1}$ and $t_{ij}\ne t_{i+1,j}$ whenever both
indices in the comparison exist. Thus $t_{ij}$ is the ordinary
main-diagonal trace, not the sum of all four entries. Rows and columns are
labeled, the boundary is not periodic, and no symmetry quotient is taken.
Write
\[
 F_w(z)=\sum_{n\ge1}a_w(n)z^n.
\]
Its constant coefficient is zero by definition. This convention does not
assign zero to the number $2^w$ of single-row arrays.

\begin{theorem}\label{thm:main}
For $2\le w\le6$, let $c_w=(c_{w,1},\ldots,c_{w,d_w})$ and $N_w$ be
as follows:
\begin{center}
\begin{tabular}{crrrl}
\toprule
$w$&OEIS&$d_w$&$N_w$&$c_w$\\
\midrule
2&A185761&2&3&$(1,4)$\\
3&A185762&3&4&$(1,6,1)$\\
4&A185763&5&7&$(1,10,3,-5,1)$\\
5&A185764&7&10&$(1,17,9,-22,-1,1,-6)$\\
6&A185765&9&13&$(1,30,24,-95,-15,64,-116,16,16)$\\
\bottomrule
\end{tabular}
\end{center}
Then
\begin{equation}\label{eq:rec}
 a_w(n)=\sum_{j=1}^{d_w}c_{w,j}a_w(n-j)\qquad(n\ge N_w).
\end{equation}
If $D_w(z)=1-\sum_{j=1}^{d_w}c_{w,j}z^j$, then
\begin{equation}\label{eq:gf}
 F_w(z)=\frac{8zP_w(z)}{D_w(z)},
\end{equation}
where
\begin{align*}
P_2(z)={}&2+3z,\\
P_3(z)={}&5+10z+z^2,\\
P_4(z)={}&13+33z-16z^3+7z^4-z^5,\\
P_5(z)={}&33+108z-43z^2-249z^3-69z^4+5z^5\\
 &\hspace{12mm}+43z^6+81z^7+18z^8,\\
P_6(z)={}&85+347z-483z^2-2718z^3-1375z^4+2104z^5\\
 &\hspace{12mm}+6276z^6+4796z^7-7734z^8-2893z^9\\
 &\hspace{12mm}+3400z^{10}-520z^{11}.
\end{align*}
The four functions $F_2,\ldots,F_5$ are Barker's formulas; the recurrence for
$w=6$ is Hardin's empirical formula in A185765~\cite{oeis6}.
\end{theorem}

\begin{corollary}\label{cor:siblings}
Hardin's remaining three empirical recurrences in this family,
A185766--A185768~\cite{oeis7,oeis8,oeis9}, also hold. In the notation
of~\eqref{eq:rec}, their coefficients and thresholds are
\begin{align}
c_7={}&(1,54,60,-400,-155,924,-1220,-76,1725,\notag\\
 &\qquad -642,-412,-88,100,40), &N_7&=19,\notag\\
c_8={}&(1,98,145,-1606,-1043,9813,-11042,-14222,67346,\notag\\
 &\qquad -40216,-64047,36082,29210,8981,-23650,-12520,\notag\\
 &\qquad 7365,448,-196), &N_8&=25,\label{eq:siblingcoeff}\\
c_9={}&(1,179,341,-6316,-6138,92272,-87872,-479803,\notag\\
 &\qquad 1815915,-852208,-6873004,8175787,6036792,\notag\\
 &\qquad -5492010,-12389568,3504328,38808492,-4735254,\notag\\
 &\qquad -36745560,-36433267,12744011,28480781,9986028,\notag\\
 &\qquad -2174189,-3904760,-845298,-118044,-60864,-3296,-32),
 &N_9&=37.\notag
\end{align}
Here $d_7=14$, $d_8=19$ and $d_9=30$.
\end{corollary}
\begin{proof}
Use the same array/path bijection and propagation in
Lemma~\ref{lem:paths} below, with the continuation identities and boundary
indices in Lemma~\ref{lem:finite}. They give precisely the printed
thresholds $n>18$, $n>24$, and $n>36$.
\end{proof}

\section{Transfer paths and a finite certificate}
Let $R_w=\{0,1\}^w$. For rows $r,s\in R_w$ define the trace vector
\[
 \tau(r,s)_j=r_j+s_{j+1}\quad(0\le j<w-1).
\]
Let $S_w$ consist of the ordered pairs $(r,s)$ for which consecutive
coordinates of $\tau(r,s)$ differ. Put an edge from $(r,s)$ to $(s,t)$
precisely when $(s,t)\in S_w$ and
\begin{equation}\label{eq:edge}
 \tau(r,s)_j\ne\tau(s,t)_j\quad(0\le j<w-1).
\end{equation}
There is at most one such edge for each ordered pair of states. Let $T_w$
be the adjacency matrix, with the source state indexing its row. The
column vector $\one$ has one in every coordinate, and
$q_k=T_w^k\one$ for $k\ge0$.

\begin{lemma}\label{lem:paths}
For $n\ge1$,
\begin{equation}\label{eq:count}
 a_w(n)=\one^{\mathsf T}T_w^{n-1}\one.
\end{equation}
If, for nonnegative $b$ and coefficients $c_1,\ldots,c_d$,
\begin{equation}\label{eq:identity}
 q_{b+d}=\sum_{j=1}^d c_jq_{b+d-j},
\end{equation}
then $a_w(n)=\sum_{j=1}^dc_ja_w(n-j)$ for every $n\ge b+d+1$.
\end{lemma}
\begin{proof}
An array with successive rows $r_0,\ldots,r_n$ determines the path
$(r_0,r_1),\ldots,(r_{n-1},r_n)$. Membership in $S_w$ imposes exactly the
horizontal inequalities, and each edge imposes exactly the vertical
inequalities. Conversely a path recovers every row uniquely, because
successive states share their middle row. This is a bijection, including
the case $n=1$ of a path with no edges, and proves~\eqref{eq:count}.

For $n\ge b+d+1$, multiply~\eqref{eq:identity} on the left by
$\one^{\mathsf T}T_w^{n-1-b-d}$. Every resulting exponent is nonnegative,
and~\eqref{eq:count} gives the asserted recurrence.
\end{proof}

\begin{lemma}\label{lem:finite}
For the coefficients of Theorem~\ref{thm:main} and
Corollary~\ref{cor:siblings}, identity~\eqref{eq:identity}
holds with the following shifts. In addition, the scalar recurrence holds
at each of the indicated boundary indices.
\begin{center}
\begin{tabular}{rrrl}
\toprule
$w$&$b$&$b+d_w$&boundary indices\\
\midrule
2&1&3&3\\
3&2&5&4,5\\
4&3&8&7,8\\
5&3&10&10\\
6&5&14&13,14\\
7&6&20&19,20\\
8&7&26&25,26\\
9&8&38&37,38\\
\bottomrule
\end{tabular}
\end{center}
\end{lemma}
\begin{proof}
This is a finite integer calculation. For each tuple
$(w,c_w,b,N_w)$ in the tables above, the specification below
enumerates every binary row, every admissible row pair, and every possible
next row; it computes $q_0,\ldots,q_{b+d_w}$ and tests each coordinate of
the identity. Finally it tests exactly the boundary scalar equalities
listed in the table. It uses neither fitted coefficients nor a cutoff on
reachable states. The accompanying \texttt{verify.py} additionally
compares counts with direct enumeration of the cells of small arrays.

\begin{Verbatim}[fontsize=\small]
from itertools import product

def check(width, coeff, shift, start):
    rows = list(product((0,1), repeat=width))
    trace = {(upper,lower): tuple(
        upper[col] + lower[col+1] for col in range(width-1))
        for upper in rows for lower in rows}
    states = [pair for pair, values in trace.items()
        if all(left != right
               for left,right in zip(values,values[1:]))]
    index = {pair: pos for pos,pair in enumerate(states)}
    edges = [[index[middle,lower] for lower in rows
        if (middle,lower) in index and all(left != right
            for left,right in zip(trace[upper,middle],
                                  trace[middle,lower]))]
        for upper,middle in states]
    degree = len(coeff)
    vectors = [[1]*len(states)]
    for exponent in range(shift+degree):
        vectors.append([sum(vectors[-1][target]
            for target in successors) for successors in edges])
    assert all(vectors[-1][pos] == sum(
        coeff[lag-1]*vectors[-1-lag][pos]
        for lag in range(1,degree+1))
        for pos in range(len(states)))
    counts = [0] + [sum(vector) for vector in vectors]
    for number in range(start,shift+degree+1):
        assert counts[number] == sum(
            coeff[lag-1]*counts[number-lag]
            for lag in range(1,degree+1))
\end{Verbatim}
\end{proof}

\begin{proof}[Proof of Theorem~\ref{thm:main}]
Lemmas~\ref{lem:paths} and~\ref{lem:finite} prove~\eqref{eq:rec}, including
the boundary indices. It remains to identify the numerators in~\eqref{eq:gf}.
For this purpose set $\widetilde a_w(0)=0$ and
$\widetilde a_w(n)=a_w(n)$ for $n\ge1$. The coefficient of $z^m$ in
$D_wF_w$ is
\begin{equation}\label{eq:num}
 \widetilde a_w(m)-\sum_{j=1}^{\min(d_w,m)}
 c_{w,j}\widetilde a_w(m-j).
\end{equation}
It vanishes for $m\ge N_w$, by the recurrence. The values for $m<N_w$
come from~\eqref{eq:count}; substitution in~\eqref{eq:num} gives the
displayed $8zP_w$. For example, at width four the coefficients for
$m=1,\ldots,6$ are $(104,264,0,-128,56,-8)$, so
$P_4=13+33z-16z^3+7z^4-z^5$.
The same exact coefficient calculation for all five widths is included
in \texttt{verify.py}. Finally $D_w(0)=1$, so $D_w$ is invertible as a
formal power series and~\eqref{eq:gf} follows.
\end{proof}

\section{The closed form and the parent table}
\begin{corollary}\label{cor:closed}
Put $s=\sqrt{17}$ and $r_\pm=(1\pm s)/2$. For every $n\ge1$,
\begin{equation}\label{eq:closed}
 a_2(n)=\frac{(13+3s)r_+^n+(-13+3s)r_-^n}{s}.
\end{equation}
This is Barker's closed form in A185761.
\end{corollary}
\begin{proof}
Both roots satisfy $r_\pm^2=r_\pm+4$, so the right-hand side satisfies
the width-two recurrence. It has values $16$ and $40$ at $n=1,2$,
respectively. These are the coefficients of $F_2$ at $z,z^2$; induction
using~\eqref{eq:rec} proves the formula. Its range begins at $n=1$;
it is not a formula for the constant coefficient of $F_2$.
\end{proof}

To identify the earlier width-two result precisely, Riedel~\cite{riedel}
counts words of length $n$ over $\{0,1,2,3\}$ avoiding
$00,11,22,33,03,30$. If $b(n)$ is this count, including $b(0)=1$, he proves
\[
 H(z)=\sum_{n\ge0}b(n)z^n=\frac{1+3z+2z^2}{1-z-4z^2}.
\]
For an array of width two, encode each diagonal pair
$(A_{i,0},A_{i+1,1})$, for $0\le i<n$, by
$01\mapsto0$, $10\mapsto3$, $00\mapsto1$, and $11\mapsto2$.
These pairs use disjoint cells. Equal successive traces are exactly the
six forbidden word pairs: the trace-one class contributes
$00,03,30,33$, and the trace-zero and trace-two classes contribute
$11$ and $22$. The remaining two cells are independently free, so
$a_2(n)=4b(n)$ for $n\ge1$. Consequently
\[
 F_2(z)=4\bigl(H(z)-1\bigr)=\frac{8z(2+3z)}{1-z-4z^2}.
\]
Thus Riedel's 23 December 2016 argument is a prior proof of the equivalent
enumeration; the transfer calculation above gives a uniform alternative.

A later antecedent is Howroyd's 9 January 2025 counting argument in
A183774~\cite{howroyd}. Write $h(k)=\mathrm{A183774}(k)$: this is half
the number of $(k+1)\times2$ binary arrays with no contiguous block
containing exactly two ones. His bottom-row classification gives
$h(k)=h(k-1)+4h(k-2)$, with $h(0)=2$ and $h(1)=5$.
There is a direct connection to the present arrays. Let
$p_i=(A_{i,0},A_{i+1,1})$, and let row $i$ of a new $n\times2$ array
be $p_i$ for even $i$ and its bitwise complement for odd $i$.
The block on rows $i,i+1$ then contains
$2+(-1)^i(t_i-t_{i+1})$ ones. Its total differs from two precisely when
the original adjacent traces differ. The construction is reversible
apart from the two free corners, giving
$a_2(n)=8h(n-1)$ for $n\ge1$. Thus Howroyd's recurrence also implies
the width-two formulas; the explicit identification of the two array
models is the deduction given here.

There is also a small hand calculation for the width-two transfer matrix.
Partition its sixteen states by $(q_1,q_2)$, whose possible values are
$(2,4),(2,6),(4,10)$. These classes have sizes $4,8,4$ and quotient matrix
\[
 M=\begin{pmatrix}0&2&0\\1&0&1\\1&2&1\end{pmatrix}.
\]
The three continuation vectors are
$M\one=(2,2,4)^{\mathsf T}$,
$M^2\one=(4,6,10)^{\mathsf T}$, and
$M^3\one=(12,14,26)^{\mathsf T}=M^2\one+4M\one$.
Weighting by $(4,8,4)$ gives $16,40,104$ at $n=1,2,3$ and hence the
same recurrence and closed form.

\begin{corollary}\label{cor:parent}
For the table $T(n,k)$ of A185769~\cite{oeistable},
\[
 T(n,k)=T(k,n)\quad(n,k\ge1),\qquad T(n,w-1)=a_w(n).
\]
Thus Theorem~\ref{thm:main} and Corollary~\ref{cor:siblings} determine
the first eight columns and, by transposition, the first eight rows.
For every fixed width $w\ge2$, $F_w$
is rational.
\end{corollary}
\begin{proof}
Transposition sends an $(n+1)\times(k+1)$ binary array to a
$(k+1)\times(n+1)$ array. It preserves each block trace and exchanges
horizontal and vertical adjacency, so it is a bijection of the counted
sets. The column identity is their definition.
For fixed $w$, the state set $S_w$ is finite, and~\eqref{eq:count} gives
\[
 F_w(z)=z\one^{\mathsf T}(I-zT_w)^{-1}\one.
\]
The inverse exists formally because the constant matrix is $I$; its
adjugate expression has polynomial numerator and denominator with
denominator constant term one. Therefore $F_w$ is rational.
\end{proof}

For widths seven through nine an explicit numerator follows without
another recurrence calculation. Set $\widetilde a_w(0)=0$ and use the
transfer formula for $a_w(1),\ldots,a_w(N_w-1)$. Then
\[
 F_w(z)=\frac{Q_w(z)}{1-\sum_{j=1}^{d_w}c_{w,j}z^j},\qquad
 Q_w(z)=\sum_{m=1}^{N_w-1}\left(\widetilde a_w(m)-
 \sum_{j=1}^{\min(d_w,m)}c_{w,j}\widetilde a_w(m-j)\right)z^m.
\]
Indeed, the coefficient calculation~\eqref{eq:num} applies unchanged,
and all coefficients from degree $N_w$ onward vanish.
The supporting integer certificates give these three polynomials in full.

\begin{corollary}\label{cor:divisibility}
Every $a_w(n)$, for $n\ge1$ and $w\ge2$, is divisible by eight.
\end{corollary}
\begin{proof}
The upper-right and lower-left cells of the array occur in no trace:
the former cannot be an upper-left or lower-right cell of a block,
and neither can the latter. They are distinct and independently free.
On the remaining cells simultaneous complementation sends every trace
$t$ to $2-t$, preserving all inequalities. This is a fixed-point-free
involution on the admissible assignments to those cells, since at least
two cells remain and a binary value differs from its complement.
The two free bits and this
involution partition the arrays into classes of size eight.
\end{proof}

\section{Coverage and supporting sources}
The following table distinguishes the assertions in the posing entries
from consequences that those entries do not pose as conjectures.
\begin{center}
\begin{tabular}{p{0.24\linewidth}p{0.58\linewidth}}
\toprule
Source item&Coverage\\
\midrule
A185761, formula fields&Prior proofs cited: Riedel's equivalent word
enumeration (2016), followed by Howroyd's related array argument (2025),
imply the recurrence, Barker's generating function and closed form through
the explicit encodings above; a uniform alternative proof is given here.\\
A185762--A185764, formula fields&Proved: all three recurrences and
Barker's three generating functions, with the stated starting indices.\\
A185765, formula field&Proved: width-six recurrence for $n>12$. Its
generating function is not posed there; it is proved here.\\
A185766--A185768, formula fields&Proved: all three empirical recurrences,
with thresholds $n>18$, $n>24$, and $n>36$, respectively.\\
A185769, name and table&The first eight row/column recurrences follow
above. Transposition symmetry, fixed-width rationality and divisibility
by eight are not posed there; they are proved above. Specific low-order
formulas at widths at least ten are not covered.\\
\bottomrule
\end{tabular}
\end{center}

The written argument supplies the array/path bijection, propagation of
the finite identities, generating-function deduction, closed form, and
transposition corollary. The integer enumeration supplies the explicit
continuation and boundary identities in Lemma~\ref{lem:finite}.
The following pinned Lean declarations have corresponding statement scopes:
\begin{itemize}
\item \href{\leanbase HardinA185761Recurrence.lean}{\texttt{HardinA185761Recurrence.result}}:
the literal width-two array count, its two initial values, recurrence,
generating function, and real closed form.
\item \href{\leanbase HardinA185762Recurrence.lean}{\texttt{HardinA185762Recurrence.result}}:
the literal width-three count, its three initial values, recurrence, and
generating function.
\item \href{\leanbase HardinA185763Recurrence.lean}{\texttt{HardinA185763Recurrence.result}}:
the literal width-four count, its six initial values, recurrence for
$n\ge7$, and generating function.
\item \href{\leanbase HardinA185764Recurrence.lean}{\texttt{HardinA185764Recurrence.result}}:
the literal width-five count, its nine initial values, recurrence for
$n\ge10$, and generating function.
\item \href{\leanbase HardinA185765Recurrence.lean}{\texttt{HardinA185765Recurrence.result}}:
the literal width-six count's recurrence for $n\ge13$, using
\href{\leanbase HardinA185765Transfer.lean}{\texttt{HardinA185765Transfer}}.
The displayed width-six generating function is a written consequence
above.
\end{itemize}
All five result declarations are in namespace
\texttt{D5.S3.Combinatorics.TraceAvoidingArrays} followed by the displayed
module name. The shared
\href{\leanbase TransferMatrix.lean}{\texttt{TransferMatrix.cardMatrix}}
gives the finite local-constraint counting identity used by the first
four modules.
Corollary~\ref{cor:siblings} uses the written propagation proof and the
finite integer identities specified here; the linked declarations cover
widths two through six, not the three additional recurrences.
The transposition, rationality and divisibility corollaries have written
proofs above.

\begingroup\small
\begin{thebibliography}{12}
\bibitem{riedel} Marko Riedel, answer to ``Recursive combinatorial solution
for number of sequences of $\{0,1,2,3\}$,'' Mathematics Stack Exchange,
\href{https://math.stackexchange.com/a/2070074}{answer 2070074},
23 December 2016; \href{https://math.stackexchange.com/revisions/2070074/3}{revision 3},
23:49 UTC, contains the four-state and two-state counting equations.
\bibitem{howroyd} Andrew Howroyd, comment of 9 January 2025 in R.\ H.\
Hardin's OEIS \href{https://oeis.org/A183774/internal}{A183774},
comment field beginning ``Let b(n) and c(n)'': bottom-row counting
equations. The entry's offset is zero.
\bibitem{stanley} Richard P.\ Stanley, \emph{Enumerative Combinatorics},
Volume 1, second-edition author version of 15 July 2011,
\href{https://math.mit.edu/~rstan/ec/ec1.pdf}{Section 4.7, pp.\ 573--574},
Theorems 4.7.1 and 4.7.2.
\bibitem{oeis2} R.\ H.\ Hardin, OEIS \href{https://oeis.org/A185761}{A185761},
name and formula fields, revision 8, 22 February 2018.
Sequence entered 3 February 2011; generating function and closed-form
conjectures by Colin Barker, 22 February 2018.
\bibitem{oeis3} R.\ H.\ Hardin, OEIS \href{https://oeis.org/A185762}{A185762},
name and formula fields, revision 9, 16 April 2018.
Sequence entered 3 February 2011; generating function by Colin Barker,
16 April 2018.
\bibitem{oeis4} R.\ H.\ Hardin, OEIS \href{https://oeis.org/A185763}{A185763},
name and formula fields, revision 9, 16 April 2018.
Sequence entered 3 February 2011; generating function by Colin Barker,
16 April 2018.
\bibitem{oeis5} R.\ H.\ Hardin, OEIS \href{https://oeis.org/A185764}{A185764},
name and formula fields, revision 10, 16 April 2018.
Sequence entered 3 February 2011; generating function by Colin Barker,
16 April 2018.
\bibitem{oeis6} R.\ H.\ Hardin, OEIS \href{https://oeis.org/A185765}{A185765},
name and formula fields, revision 9, 3 October 2025.
Sequence entered 3 February 2011.
\bibitem{oeis7} R.\ H.\ Hardin, OEIS \href{https://oeis.org/A185766}{A185766},
name and formula fields, revision 9, 3 October 2025.
Sequence entered 3 February 2011; empirical recurrence in formula field.
\bibitem{oeis8} R.\ H.\ Hardin, OEIS \href{https://oeis.org/A185767}{A185767},
name and formula fields, revision 9, 3 October 2025.
Sequence entered 3 February 2011; empirical recurrence in formula field.
\bibitem{oeis9} R.\ H.\ Hardin, OEIS \href{https://oeis.org/A185768}{A185768},
name and formula fields, revision 9, 3 October 2025.
Sequence entered 3 February 2011; empirical recurrence in formula field.
\bibitem{oeistable} R.\ H.\ Hardin, OEIS
\href{https://oeis.org/A185769}{A185769}, name and table, revision 9,
3 October 2025. Sequence entered 3 February 2011.
\end{thebibliography}
\noindent\small Copyright 2026 Haobo Ma and Wenlin Zhang.
Licensed under \href{https://creativecommons.org/licenses/by/4.0/}{Creative
Commons Attribution 4.0 International (CC BY 4.0)}.
\endgroup
\end{document}
